\documentclass{article}
\usepackage{amsmath}
\usepackage{amsthm}
\usepackage{amsfonts}
\usepackage{enumitem}

\newtheorem{theorem}{Theorem}

\author{Catherine Reid}
\title{Metric Spaces}

\begin{document}

\maketitle

\textbf{Exercise 3.1:} Let $X$ be the set of continuous functions $f:[a, b] \to \mathbb{R}$. Let $d^*$ be the distance function
on $X$ defined by
\begin{align*}
    d^*(f, g) = \int_{b}^{a}  \left|f(t) - g(t)\right|\,dt,
\end{align*}
for $f, g \in X$. For each $f \in X$, set
\begin{align*}
    I(f) = \int_{b}^{a} f(t)\,dt
\end{align*}
Prove that the function $I:(X, d^*) \to (\mathbb{R}, d)$ is continuous.

\begin{proof}
    Let $g \in X$. Let $\varepsilon > 0$. Choose $\delta = \varepsilon$. Suppose that $f \in X$ is such
    that $d^*(f, g) < \delta$, that is
    \begin{align*}
        \int_{b}^{a}  \left|f(t) - g(t)\right|\,dt < \delta
    \end{align*}
    Next we note that
    \begin{align*}
        d(I(f), I(g)) &= \left|\int_{b}^{a} f(t)\,dt - \int_{b}^{a} g(t)\,dt\right|\\
        &= \left|\int_{b}^{a} f(t) - g(t)\,dt\right|\\
        &\le \int_{b}^{a} \left|f(t) - g(t)\right|\,dt\\
        &< \delta\\
        &= \varepsilon
    \end{align*}
    
    Therefore, given $\varepsilon > 0$, there is a $\delta > 0$ such that $d(I(f), I(g)) < \varepsilon$ whenever
    $d^*(f, g) < \delta$.
\end{proof}

\pagebreak
\textbf{Exercise 3.2:}  Let $(X_i, d_i)$, $(Y_i, d^\prime_i)$, $i = 1, \dots, n$ be metric spaces. Let $f_i: X_i \to Y_i$,
$i = 1, \dots, n$ be continuous functions. Let
\begin{align*}
    X = \prod_{i = 1}^{n} X_i & & \text{and} & & Y = \prod_{i = 1}^{n} Y_i
\end{align*}
and convert $X$ and $Y$ into metric spaces in the standard manner. Define the function $F: X \to Y$ by
\begin{align*}
    F(x_1, \dots, x_n) = (f_1(x_1), \dots, f_n(x_n))
\end{align*}
Prove that $F$ is continuous

\begin{proof}
    Let $a = (a_1, \dots, a_n) \in X$. Let $\varepsilon > 0$. Since each $f_i$ is continuous at $a_i$, we know that
    given $\varepsilon > 0$, there exists a $\delta_i > 0$ such that $x_i \in X_i$ and $d_i(a_i, x_i) < \delta_i$ imply that
    $d^\prime_i(f_i(x_i), f_i(a_i)) < \varepsilon$. Choose $\delta = \min_{1 \le i \le n} \delta_i$.

    Now, suppose that $x = (x_i, \dots, x_n) \in X$ is such that $d(x, a) < \delta$. That is
    \begin{align*}
        \max_{1 \le i \le n} \left\{d_i(x_i, a_i)\right\} < \delta
    \end{align*}

    Then
    \begin{align*}
        d^\prime(F(x), F(a)) &= d^\prime((f_1(x_1), \dots, f_n(x_n)), (f_1(a_1), \dots, f_n(x_n)))\\
        &= \max_{1 \le i \le n} \left\{d^\prime_i(f_i(x_i), f_i(a_i))\right\}\\
        &< \varepsilon
    \end{align*}
    We justify the final conclusion by pointing out that since every $d_i(x_i, a_i) < \delta$, it is the case
    that every $(d^\prime_i(f_i(x_i), f_i(a_i))) < \varepsilon$.
\end{proof}

\pagebreak
\textbf{Theorem 4.6:} Let $f(X, d) \to (Y, d^\prime)$. $f$ is continuous at a point $a \in X$ if and only if for each neighbourhood $M$ of $f(a)$
there is a corresponding neighbourhood $N$ of $a$ such that
\begin{align*}
    f(N) \subset M,
\end{align*}
or, equivalently
\begin{align*}
    N \subset f^{-1}(M)
\end{align*}

\begin{proof}
    First suppose that $f$ is continuous at the point $a \in X$. For a neighbourhood $M$ of $f(a)$ we have an $\varepsilon > 0$ such
    that $B(f(a);\,\varepsilon) \subset M$. Since $f$ is continuous at $a$, there is a $\delta > 0$ such that $f(B(a;\,\delta)) \subset B(f(a);\,\varepsilon)$.
    $N = B(a;\,\delta)$ is a neighbourhood of $a$. Therefore
    \begin{align*}
        f(N) = f(B(a;\,\delta)) \subset B(f(a);\,\varepsilon) \subset M
    \end{align*}
    
    Conversely, suppose that $f$ satisfies the condition that for each neighbourhood of $M$ of $f(a)$, there is a corresponding neighbourhood
    $N$ of $a$, such that $f(N) \subset M$. Let $\varepsilon > 0$ be given. We need to show that there exists a $\delta > 0$ such that
    \begin{align*}
        f(B(a;\,\delta)) \subset B(f(a);\,\varepsilon)
    \end{align*}
    Now, $B(f(a);\,\varepsilon) = M$ is a neighbourhood of $f(a)$ and therefore there is a neighbourhood $N$ of $a$ such that $f(N) \subset M$.
    Since $N$ is a neighbourhood of $a$, there is a $\delta > 0$ such that $B(a;\,\delta) \subset N$. Therefore
    \begin{align*}
        f(B(a;\,\delta)) \subset f(N) \subset M = B(f(a);\,\varepsilon  )
    \end{align*}
\end{proof}
\pagebreak

\textbf{Exercise 4.1:} Let $(X, d)$ be a metric space such that $d(x, y) = 1$ whenever $x \ne y$. Let $a \in X$. Prove that $\{a\}$ is a
neighbourhood of $a$ and constitutes a basis for the system of neighbourhoods at $a$. Prove that every subset of $X$ is a neighbourhood
of each of its points.

\begin{proof}
    Let $\delta = 1/2$. Then, since $d(a, x) = 1$ when $a \ne x$, $B(a;\,\delta) = \{a\}$. By definition,
    $\{a\}$ is therefore a neighbourhood.

    Consider a neighbourhood $M$ of $a$. By Theorem 4.8, we know that $M$ contains $a$. Therefore $\{a\} \subset M$, and 
    $\{a\}$ constitutes a basis for the system of neighbourhoods at $a$.

    Now, consider $Y$ where $Y \subset X$. For any $y \in Y$, $\{y\}$ is a neighbourhood of $y$. By Theorem 4.8 this means $Y$
    is also a neighbourhood of $y$. Therefore any subset of $X$ is a neighbourhood of each of its points
\end{proof}

\textbf{Exercise 4.2:} Let $a \in \mathbb{R}$ and $f: \mathbb{R} \to \mathbb{R}$ be defined by
\begin{align*}
    f(x) = \begin{cases} 
      0 & x \le a \\
      1 & x > a
   \end{cases}
\end{align*}
Prove that $f$ is not continuous at $a$, but is continuous everywhere else
\begin{proof}
    Firstly we will show that $f$ is not continuous at $a$. To do this we must show that for some $\varepsilon_0 > 0$ and any $\delta > 0$,
    we have that there is some $y \in \mathbb{R}$, such that $y \in f(B(a;\,\delta))$ and $y \notin B(f(a);\,\varepsilon)$.

    Let $\varepsilon_0 = \frac{1}{2}$ and let $\delta > 0$. Choose $x = a + \frac{\delta}{2}$. Now, $f(x) = 1$, so $1 \in f(B(a;\,\delta))$.
    But, $f(a) = 0$ and $\left|0 - 1\right| \ge \frac{1}{2}$, and so $1 \notin B(f(a);\,\varepsilon_0)$. Therefore $f(a)$ is not continuous at $a$.

    For any point $x$ where $x \ne a$, we have that $d(x, a) = \delta$ for some $\delta > 0$. Therefore $a \notin B(x;\,\delta)$,
    and so due to the definition of $f$, $f(B(a;\,\delta))$ is constant. This implies that $f$ is continuous at $x$, and so is
    continuous at any point other than $a$.
\end{proof}

\textbf{Exercise 4.3:} Let $f: X \to Y$ be a function into a metric space $Y$. Let $a \in X$ at let $\mathfrak{B}_{f(a)}$ be a basis for
the neighbourhood system at $f(a)$. Prove that $f$ is continuous at $a$ if and only if for each $N \in \mathfrak{B}_{f(a)},\,f^{-1}(N)$ is
a neighbourhood for $a$.

\begin{proof}
    First we will assume $f$ is continuous at $a$. Each $N \in \mathfrak{B}_{f(a)}$ is a neighbourhood of $f(a)$. Therefore by Theorem 4.6 we know that there exists a neighbourhood of $M$ of $a$ such that
    $M \subset f^{-1}(N)$. Thus, by Theorem 4.8, we can conclude that $f^{-1}(N)$ is also a neighbourhood of $a$.

    Conversely, assume that for each $N \in \mathfrak{B}_{f(a)},\,f^{-1}(N)$ is a neighbourhood of $a$. Now, consider a neighbourhood $M$ of
    $f(a)$. By definition for some $N \in \mathfrak{B}_{f(a)}$, we have $N \subset M$, so $f^{-1}(N) \subset f^{-1}(M)$. Since, $f^{-1}(N)$
    is a neighbourhood of $a$, by Theorem 4.8 we can conclude that $f(a)$ is continuous
\end{proof}

\textbf{Exercise 4.4:} Let $a$ be a point on the real line $\mathbb{R}$. Prove that each of the following subsets of collections of 
subsets of $\mathbb{R}$ constitutes a basis for the system of neighbourhoods at $a$.

\begin{enumerate}[label=\roman*)]
    \item All the closed intervals of the form $[a - \varepsilon,\,a + \varepsilon],\, \varepsilon > 0$.
    
    \begin{proof}
        Consider a neighbourhood $N$ of $a$. By definiton there is a $\delta > 0$ such that
        \begin{align*}
            (a - \delta,\,a + \delta) \subset N
        \end{align*}

        Let $\varepsilon_0 = \delta / 2$. Then we have
        \begin{align*}
            [a - \varepsilon_0,\,a + \varepsilon_0] \subset (a - \delta,\,a + \delta) \subset N
        \end{align*}
        Thus an interval of the form $[a - \varepsilon,\,a + \varepsilon],\, \varepsilon > 0$ is a subset of every neighbourhood $N$ of $a$,
        so the collection is a basis for the system of neighbourhoods of $a$
    \end{proof}

    \item All open balls $B(a;\, \varepsilon)$, $\varepsilon$ a positive rational number.
    
    \begin{proof}
        Consider a neighbourhood $N$ of $a$. By definiton there is a $\delta > 0$ such that
        \begin{align*}
            B(a;\,\delta) \subset N
        \end{align*}

        By the density of the rationals in the reals, there exists a rational number $\varepsilon_0$ such that $0 < \varepsilon_0 < \delta$.
        Therefore we have
        \begin{align*}
            B(a;\,\varepsilon_0) \subset B(a;\,\delta) \subset N
        \end{align*}
        Thus an open ball $B(a;\, \varepsilon)$ where $\varepsilon$ is a positive rational number is a subset of every neighbourhood $N$ of $a$,
        so the collection is a basis for the system of neighbourhoods of $a$.
    \end{proof}

    \item All open balls $B\left(a;\,\dfrac{1}{n}\right)$, $n$ is a positive integer.
    
    \begin{proof}
        Consider a neighbourhood $N$ of $a$. By definiton there is a $\delta > 0$ such that
        \begin{align*}
            B(a;\,\delta) \subset N
        \end{align*}
        By the Archimedian Property of the reals, there exists a natural number $n_0$ such that $\dfrac{1}{n_0} < \delta$. Therefore we have
        \begin{align*}
            B\left(a;\,\dfrac{1}{n_0}\right) \subset B(a;\,\delta) \subset N
        \end{align*}
        Thus an open ball $B\left(a;\,\dfrac{1}{n}\right)$ where $n$ is a positive integer is a subset of every neighbourhood $N$ of $a$,
        so the collection is a basis for the system of neighbourhoods of $a$. 
    \end{proof}
\end{enumerate}

Finally, show that no finite collection of subsets of $\mathbb{R}$ can be a basis for the system of neighbourhoods of $a$.

\begin{proof}
    Consider for the sake of contradiction that a finite collection of subsets of $\mathbb{R}$ can be the basis the the system of
    neighbourhoods of $a$. Call this collection $\mathfrak{B}_a$.
    
    Now, for each element $B_i \in \mathfrak{B}_a$, there exists a $\delta_i > 0$ such that $B(a;\,\delta_i) \subset B_i$.
    Since $\mathfrak{B}_a$ is finite, there exists an ordering such that
    \begin{align*}
        \delta_1 \le \delta_2 \le \dots \le \delta_n
    \end{align*}
    That is
    \begin{align*}
        B(a;\,\delta_1) \subset B(a;\,\delta_2) \subset \dots \subset B(a;\,\delta_n)
    \end{align*}
    Now, let $\varepsilon_0 = \dfrac{\delta_1}{2}$. Therefore we have
    \begin{align*}
        B(a;\,\varepsilon_0) \subset B(a;\,\delta_1) \subset B(a;\,\delta_2) \subset \dots \subset B(a;\,\delta_n)
    \end{align*}
    By definition, $B(a;\,\varepsilon_0)$ is a neighbourhood of $a$. But we have shown that it is a subset of every element of 
    $\mathfrak{B}_a$, contradicting the fact that $\mathfrak{B}_a$ is a basis of the system of neighbourhoods of $a$. Therefore
    we can conclude that no finite collection of subsets of $\mathbb{R}$ can be the basis of the system of neighbourhoods of $a$.
\end{proof}
\pagebreak

\textbf{Exercise 4.6:} Let $a$ and $b$ be distinct points of a metric space $X$. Prove that there are neighbourhoods $N_a$ and $N_b$
of $a$ and $b$ respectively such that $N_a \cap N_b = \emptyset$

\begin{proof}
    Since $a \ne b$, $d(a, b) > 0$. Let $\varepsilon = d(a, b)$. Now, let $N_a = B(a;\,\frac{\varepsilon}{2})$ and 
    $N_b = B(b;\,\frac{\varepsilon}{2})$. Consider the point $x \in B(a;\,\frac{\varepsilon}{2})$. By definition
    $d(a, x) \le \frac{\varepsilon}{2}$. We have that
    \begin{align*}
        d(a, b) &\le d(a, x) + d(x, b)\\
        \varepsilon &\le \frac{\varepsilon}{2} + d(x, b)\\
        \frac{\varepsilon}{2} &\le d(x, b)
    \end{align*}
    Therefore we can conclude that $x \notin B(b;\,\frac{\varepsilon}{2})$, and therefore $N_a \cap N_b = \emptyset$
\end{proof}

\textbf{Exercise 4.8:} Let $\mathbb{R}$ be the real numbers and let $f: \mathbb{R} \to \mathbb{R}$ a continuous function
Suppose that for some number $a \in \mathbb{R},\,f(a) > 0$. Prove that there is a positive number $k$ and a closed interval
$F = [a - \delta,\,a + \delta]$ for some $\delta > 0$ such that $f(x) \ge k$ for all $x \in F$.

\end{document}